\section{Reducing the total arity of a tensor power}\label{sec:rounds}

The concentration estimate of Lemma~\ref{lem:local} divides a network
into a few large stations and a core of small stations. To lower the cost
of evaluating the network we must also reduce the \emph{total} arity.
Indeed, if $u_v$ is the arity at station $v$, the batched evaluation in
Lemma~\ref{lem:batch-table} costs, per tensor copy, at most a constant
times the number of stations plus
\[
       \sum_{v:\,u_v>0}u_v^p
           \le \bigl(\max_v u_v\bigr)^{p-1}\sum_v u_v .
\]
Here $1<p<2$. We can therefore afford some growth in the largest arity
if enough incident bits disappear. The construction below makes this
trade repeatedly, keeping a bound on the sum of squared arities so that
each new network again meets the concentration hypothesis.

Throughout this part $k$ is a finite field with $\omega_k<9/4$, and
logarithms are binary. A \emph{table network} over a field $K$ has bit
wires and a table $\{0,1\}^{s}\to K$ at each station of arity $s$. Its
size is the number of stations plus the sum of their arities, including
external legs. Parallel bits between two stations form an ordered
\emph{bundle}. A label names the station's table in a dictionary fixed
for the network.

\needspace{12\baselineskip}
\begin{lemma}[Compression of a tensor power]\label{lem:macro}
Let $h=\lceil\log_2 n\rceil$, $0<\delta\le 1/10$, and
$0<\beta\le b\le 1/2$. For every sufficiently small fixed $a>0$
depending only on $k,\delta,\beta$, and all sufficiently large $n$, the
following holds. Work over an extension $K/k$ of degree $O(\log h)$
large enough for Lemma~\ref{lem:local} at the scales used below.
Suppose that an input table network has maximum arity $h^b$ or less,
$(\log h)^{O(1)}$-bit table labels, size at most $m$ with
$n\le m\le n^{1.1}$, and $e\le m h^{-2a}$ external bits. Some
$d$ with $1\le d=n^{o(1)}$ admits a table network over $K$ for its
$d$th tensor power, with all $de$ external bits separately accessible,
such that
\[
 \begin{aligned}
 \text{size}&\le md\,h^{-(1-\delta)a},&
 \text{maximum arity}&\le h^{b+(1+\delta)a}.
 \end{aligned}
\]
The output labels again have $(\log h)^{O(1)}$ bits. The cutoff for $n$
may depend on the fixed exponent in the input label-length bound.
\end{lemma}

Call one application of Lemma~\ref{lem:macro} a \emph{macro}.
Before the macro the lookup bound is $O(mh^{b(p-1)})$ per copy.
After it, the bound acquires the factor
\[
       h^{-a\{(1-\delta)-(1+\delta)(p-1)\}} .
\]
Since $p<2$, choosing $\delta$ small makes the exponent negative.
Section~\ref{sec:shaving} accumulates this saving before evaluating the
tables.

\begin{table}[bp]\centering
\small
\begin{tabular}{@{}ll@{}}
\toprule
Quantity & Meaning within one summand\\
\midrule
$H,\ N_b$ & positions of the common host and blocks using it\\
$r$ & number of blocks combined at each grouping step\\
$S,\ T_0$ & initial blocks and original copies per block, with $d=ST_0$\\
$L,\ T$ & total arity budget and working arity scale\\
$D$ & concentration parameter, fixed for all rounds\\[2pt]
$X,\ Z$ & width between core stations and total frozen charge\\
$\Delta,\ \nabla$ & degree bounds for the core host and the full host\\
$B,\ u$ & maximum packing load and unit of internal width\\
\bottomrule
\end{tabular}
\caption{Parameters for the compression of a tensor power. The budgets
$L$, $X$ and $Z$ each cover all blocks of the summand.}
\label{tab:macro}
\end{table}

The macro expands a power of the original tensor as a sum of products.
Each product has many factors, called \emph{blocks}. We combine blocks
in groups of $r$ to pack stations or rearrange their wires. Each round
first concentrates the squared arities, then uses a fixed number of
wire reductions to shorten the core, and finally restores bounded
degree so the next round can start. A wire reduction, called a
\emph{beat}, consists of several contractions followed by a reset of the
core degree. The number of rounds grows with $h$. Once $D$ is fixed,
all repetition counts \emph{within} a round are constant. Table~\ref{tab:macro} collects the budgets and the quantities they measure.

Different summands can use different wirings. The bounds in
Sections~\ref{sec:rounds}--\ref{sec:full} apply to each summand.
Section~\ref{sec:assembly} represents their sum on one fixed network
and completes the proof of Lemma~\ref{lem:macro}.

\subsection{The common graph and its initial loads}
Each summand is a product of \emph{blocks}, with disjoint lists of
external bits. In a single summand its $N_b$ blocks use the same simple
graph of possible adjacencies, the \emph{host}. An occurrence is a
pair consisting of a host position and a block. It can be empty. Each
actual internal wire joins stations in the same block, and the local
port order tells which bits at its ends correspond. Thus a host edge can
carry different bundles in different blocks. We permit the host to depend
on the summand.

At a grouping step, relabel the blocks and group them $r$ at a time, with
$r$ a power of two. In each group the old blocks are multiplied and
transformed together. New positions may be indexed by an old position, a
dyadic subinterval of the $r$ indices, and a bounded number of other
choices. There are fewer than $2r$ intervals. If $H$ denotes the number
of host positions, the padding allowed in any one step will have
\begin{equation}\label{eq:round-occurrence}
 H'\le c rH, \qquad N_b'=N_b/r,
\end{equation}
where $c$ is fixed throughout all repetitions of that step. Constants here
may depend on the round design. Transformations at a single position using
a fixed number of slots can be included in the adjacent grouping step.

At the beginning of each round, every live occurrence holds one station.
The host has degree bounded by an absolute constant, and external bits are fixed at
specified ports. With $T$ a dyadic width and $L$ a total arity budget,
the charged arity of each live station of actual arity $s_v$ is
$z_v=\max(s_v,\theta T)$, where $\theta>0$ is a small fixed dyadic
constant chosen below. Our invariants are
\begin{equation}\label{eq:round-invariants}
 \sum_v z_v\le L,\qquad \sum_v z_v^2\le LT,\qquad
 \max_v z_v\le DT/2 .
\end{equation}
All three bounds cover all blocks of the summand. The first pays for
surviving bits. The second is the input to the next local concentration,
and the third lets us pack several stations without overflowing one
cell. The floor $\theta T$ also bounds how many table descriptions can
enter a pack of a given charge. Remove every scalar station and retain
its entry as a factor of the summand.

First replace the input by a graph of bounded degree. On the ports of
each station of positive arity $s$ use a balanced binary tree. The root
receives its $s$ bits as a bundle, and an identity at each split divides
one ordered bundle in two, of lengths differing by at most one. Place
an identity at each leaf and attach the original wires and external bits
at the leaves. Connecting two of the resulting leaves also handles an input
wire that was a loop. The arities on either end of a tree edge have bounded
ratio. The leaf-to-leaf connections have matching unit width. There are
$O(\log(2+h))$ levels and $O(m\log(2+h))$ total arity. All the new
tables have short labels.

Choose a dyadic $T_0\in[4h^b,8h^b]$. Stack $T_0$ copies of the network
just constructed. At a station of arity $s$ let $v$ be the largest dyadic
integer at most $T_0/s$. Divide the copies into consecutive intervals of
length $v$. Give each interval its own position carrying the product of
the corresponding tables. Its arity $sv$ belongs to $(T_0/2,T_0]$.
Comparable arities on the two ends of each old edge imply that each
interval meets only boundedly many intervals at the other end. Each edge inherits the order of the old bundle and then of the copies,
so its two endpoints use the same bit order. The stacked graph has
bounded degree and $H_0=O(m\log(2+h))$ positions, since a station of
arity $s$ gives at most $2s$ intervals. The table at an interval is
specified by the old table label, the interval length, and a bounded
number of ordered port lengths.

Take $S$ blocks of this stack and pad $H_0$ to order $m\log(2+h)$.
Include the $d$th power of the product of the removed input scalars as
a factor of the final tensor. The starting parameters are
\begin{equation}\label{eq:round-start}
 d=T_0S,\qquad T=T_0,\qquad L=L_0=CmT_0S\log(2+h).
\end{equation}
The total charged arity is at most $L_0$ for a sufficiently large
absolute $C$. Each station has arity in $(T_0/2,T_0]$, so its square
is at most $T_0$ times its arity. This gives all three bounds in
\eqref{eq:round-invariants}. The $ed$ exterior bits retain their
individual names throughout the construction.

\section{Reducing the width between unfrozen stations}
\label{subsec:core-width}

After concentration, some stations remain frozen for the rest of this
round. The others form its \emph{core}. We shall reduce the width
between core stations before restoring bounded degree and passing all
the stations to the next round. Each core occurrence carries one
tensor. Each frozen occurrence holds a list of tensors and the charge
assigned to each when it was frozen. All exterior legs are on frozen
tensors. We discard host adjacencies between frozen stations, since they
carry no wires.

Let \(X\) be the sum, over all blocks, of the widths of the core--core
wires, counting each wire once. Let \(Z\) be the total frozen charge.
The core--frozen wires, called \emph{tethers}, have total width at most
\(Z\). For packing, give each core occurrence load equal to its arity
and each frozen occurrence load equal to the sum of its charges times
a fixed dyadic factor in \((0,1]\). Write \(B\) for an upper bound on
one load. The sum of all loads satisfies \(\mathcal M\le 2X+2Z\).
Each internal width is a multiple of an integer \(u\le B\).

The host records possible contacts across all blocks. Write \(\Delta\ge3\)
for a bound on its core degree and \(\nabla\ge\Delta\) for a bound on
its full degree. A core station can meet different frozen stations in
different blocks, so even a core of degree three can have many potential
tethers. The width estimate below is uniform in their number.

\begin{lemma}[Reduction of core width]\label{lem:core-width}
Fix a bound on the incoming core degree and assume the hypothesis of
Theorem~\ref{thm:sort}. There is a fixed sequence of grouping and
contraction steps, called a \emph{beat}, with the following bounds
in each summand, for slack \(\sigma>0\).
\[
 X_{\rm out}\le \tfrac12 X+C(Z+\sigma),\qquad
 \Delta_{\rm out}\le3,\qquad Z_{\rm out}=Z.
\]
The identity retains all exterior legs separately. The maximum load,
the reciprocal bundle unit and \(\log(2\nabla)\) increase by at most
fixed factors. These factors and \(C\) are independent of the
incoming full degree.

At a step with \(H\) host positions and \(N_b\) blocks, a power of two,
choose dyadic \(r\mid N_b\). Grouping leaves \(N_b/r\) blocks and uses
at most \(crH\) positions. It suffices to take
\(r\ge c_*\), \(c_*HN_bB/r\le\sigma\) at each step and a starting
unit divisible by the product of the subdivision factors. The
position and slack constants \(c,c_*\) may depend on the degree
bounds at the step in question.
\end{lemma}

A beat consists of several wire-saving subrounds and one reset. Each
subround trades some core width for more possible core contacts. The
reset routes the surviving wires through a core of degree three. We
shall choose the number of subrounds so that their combined saving
pays for this routing.

\subsection{A wiring identity}
Place six parties in the triangular array of Figure~\ref{fig:two-layouts}.
A pair is \emph{permitted} if it lies in one row or one column. The
desired equality wires \(AB\) and \(CD\) lie in neither. We shall express
their product as a sum of networks on permitted pairs and then contract
an entire row or column partition. This leaves a choice of two layouts
for the same summand.

\begin{figure}[htb]
\centering
\begin{tikzpicture}[x=1.06cm,y=1.32cm,font=\small,
 party/.style={circle,inner sep=1pt,minimum size=4mm},
 part/.style={rounded corners=2.8pt,line width=.65pt},
 demand/.style={draw=ink,densely dashed,thick}]
\foreach \panel/\shift/\heading in {0/0/Demanded wires,1/4.2/Row partition,2/8.4/Column partition}{
 \begin{scope}[xshift=\shift cm]
 \node[font=\sffamily\small] at (1.14,1.96) {\heading};
 \ifnum\panel=1
  \draw[part,draw=sea,fill=sea!8] (-.33,1.15) rectangle (2.61,1.65);
  \draw[part,draw=sea,fill=sea!8] (-.33,.45) rectangle (1.47,.95);
  \draw[part,draw=sea,fill=sea!8] (-.33,-.25) rectangle (.33,.25);
 \fi
 \ifnum\panel=2
  \draw[part,draw=rust,fill=rust!8] (-.33,-.25) rectangle (.33,1.65);
  \draw[part,draw=rust,fill=rust!8] (.81,.45) rectangle (1.47,1.65);
  \draw[part,draw=rust,fill=rust!8] (1.95,1.15) rectangle (2.61,1.65);
 \fi
 \ifnum\panel=0
  \draw[demand,shorten <=2.6mm,shorten >=2.6mm] (0,1.4)--(1.14,.7);
  \draw[demand,shorten <=2.6mm,shorten >=2.6mm] (2.28,1.4)
    .. controls (2.64,-.14) and (1.26,-.23) .. (0,0);
 \fi
 \foreach \x/\y/\name in {0/1.4/C,1.14/1.4/M,2.28/1.4/A,0/.7/N',1.14/.7/D,0/0/B}{
  \node[party] at (\x,\y) {\(\name\)};
 }
 \end{scope}}
\end{tikzpicture}
\caption{An equality wire between \(C\) and \(M\), for example, disappears
when the rows are contracted and survives when the columns are contracted.
Every permitted wire disappears in one layout. Across all copies of the
identity in a summand, the cheaper of the two common layouts retains at
most half the replacement width.}
\label{fig:two-layouts}
\end{figure}

An \(s\)-bit wire is the equality tensor on the two \(s\)-bit indices.
For a scale-one ordinary price over \(k\), let \(m\) be the greatest
price-distance on a permitted pair. A spanning tree of the permitted
graph supplies one bit shared by all six parties, by equating its
incident bits at each vertex. Thus \(1\le5m\). Both demanded distances
are at most \(2m\). If \(d(C,D)=2m\), the paths \(C,M,D\) and \(C,N',D\)
make \(M\) and \(N'\) the same midpoint in a CAT(0) extension. If also
\(d(A,B)=2m\), then \(C\) and \(M=N'\) are both midpoints of \(A,B\).
This would give \(C=M\) although \(d(C,M)=m>0\). Compactness of
scale-one prices therefore gives a uniform \(\lambda_*<2\) with
\[
 d(A,B)+d(C,D)\le2\lambda_*m .
\]

Choose \(\lambda_*<\lambda_0<2\). For fixed \(\eta'>0\) and arbitrarily
large \(b\), we claim that the two width-\(b\) demands admit an ordinary
sum of at most \(2^{\eta'b}\) networks, each with total permitted width
at most \(2\lambda_0b\). Indeed, choose \(h_0\) large enough that
\(\lfloor2\lambda_0h_0\rfloor>2\lambda_*h_0\). As supply, take the
direct sum over permitted pairs of a single wire on that pair of width
\(\lfloor2\lambda_0h_0\rfloor\). At each spectral point the price of
this supply exceeds that of the two width-\(h_0\) demands. Strict
spectral comparison makes the width-\(b\) demands, for \(b=h_0r\),
a restriction of boundedly many direct copies of the \(r\)th supply power.
Expanding the direct sums and applying the local restriction maps gives
the claimed networks. With \(E_{\rm perm}\) the set of permitted pairs,
a summand index takes \(r\log_2|E_{\rm perm}|+O_{h_0}(1)\) bits.
Choose \(h_0\) and then \(r\) large to fit \(\eta'b\). The identity
over \(k\) extends to every field \(K/k\). For the subround we need a
fixed number of choices even when each wire carries many copies of a bit.

\begin{lemma}[A fixed choice of widths]\label{lem:powered-gadget}
There are constants $b\in\mathbb N$ and $\lambda<2$ such that for
every $z\ge1$ the two width-$bz$ demands are the sum of a bounded
number of networks on the same six parties. Each uses at most
$2\lambda b$ bundles of $z$ bits on permitted edges. Its local
tensors are specified by $z$, the party and the index of the summand.
The index set and choices of widths (in units of $z$) can be fixed independently of $z$.
\end{lemma}

\begin{proof}
Raise the finite identity above to its $z$th tensor power. A term is
initially indexed by a sequence of $z$ selectors, with at most
$\lceil\eta'b\rceil$ bits per selector. For each permitted edge round
up the total number of bits on it to a multiple of $z$. Allowing empty groups, group the
sequences by the vector of rounded numbers in units of $z$. This
vector belongs to a fixed finite set, and rounding costs at most one
unit per permitted edge. Within each group choose a spanning tree
of the connected permitted graph and send the full sequence on each
tree edge, using at most $\lceil\eta'b\rceil$ extra units per edge.
A local table checks identical valid sequences on all incident tree
edges, membership in the group, and zeros in padding positions.
It uses the sequence to split positions on each original edge into
the prescribed slices and evaluates the product of the $z$ local
tables. Connectivity makes the encodings of a sequence unique once
the sequence is specified. Summing the tree variables therefore gives
precisely the sum for that group. The total in units of $z$ is at most
$2\lambda_0 b+|E_{\rm perm}|+5\lceil\eta'b\rceil$. Choose $\eta'$
sufficiently small and then $b$ sufficiently large to bound this by
$2\lambda b$ with $\lambda<2$.
\end{proof}


\subsection{Subrounds}
In a subround we group occurrences at each position into packs. Two
outgoing units from different packs can play the roles \(AB\) and \(CD\).
We first arrange enough such pairs, with a bound on how many packs a
single contraction can merge.

\paragraph{Balanced partitions.}
Within each group of \(r\) block indices use the dyadic intervals in their common order. Give every host edge \(e\) a propagation distance
\[
 s_e=64\lceil\log_2(2x_e)\rceil,\qquad
 x_e=\begin{cases}\Delta&\text{on core edges},\\
                     \nabla&\text{on tethers}.\end{cases}
\]
At each position mark every interval of load greater than \(AB\), with \(2\le A\le r\) dyadic. Marks require splitting. Propagate each mark to its parent at the same position and to its ancestor \(s_e\) levels higher at the other end of each incident edge \(e\), whenever the ancestor exists. The leaves, called colonies, have load at most \(AB\) and length at least \(A\). Intersecting colonies at opposite ends of \(e\) differ in length by a factor at most \(2^{s_e}\), by propagation from the parent of the smaller colony.

A common permutation of the indices inside each group gives a total colony count
\begin{equation}
 Q\le C_1\left( HN_b/r+\mathcal M/(AB)\right)
 \label{eq:core-colonies}
\end{equation}
with \(C_1\) absolute. To see this, identify the indices with \(\mathbb F_r\) in a binary basis and multiply by a uniformly random nonzero element. The probability that two distinct indices fall in one depth-\(j\) interval is at most \(2^{-j}\). At a source position in one group let \(x_i\le1/A\) be the loads divided by \(AB\), and put \(x=\sum_i x_i\). An initially marked interval has at least \(1/2\) in ordered cross products of distinct normalized loads. Thus the expected number of initial marks at depth \(j\) is at most \(2x^2 2^{-j}\). The sum of \(2^{-s_e/16}\) over incident edges is at most \(1/2\) at every vertex. Summing the resulting path weights shows that at most \(2\cdot2^{s/16}\) positions are reachable with total propagation distance at most \(s\). Fix the total upward displacement \(s\), allowing steps to parents at the same position. At each reachable position, depth-\(j\) source marks then reach at most \(2^{j-s}\) distinct ancestors. The expected count there, summed over \(j\ge s\), is at most
\[
 \sum_{j\ge s}\min\{2^{j-s},2x^2 2^{-j}\}=O(x2^{-s/2}).
\]
Summing over distances, sources and groups bounds the number of marks in expectation by \(O(\mathcal M/(AB))\). Each binary partition has one more leaf than marked internal nodes. Choose a multiplier minimizing the actual count.

\paragraph{A wire-saving subround.}
Orient all core edges. Color each block thick or thin, using one affine
functional over \(\mathbb F_2\) on its within-group index. Every position
in a block receives the same color. To pack tensors at one position,
take their tensor product and retain their wires. Any loops created by
later mergers are contracted into the local table. At each core colony gather all thick tensors in one pack. In order, close a thin pack once it reaches load \(B\). There are \(O(A)\) thin packs, each of load at most \(2B\). At a frozen colony concatenate the lists. Let
\(P=2^{64\lceil\log_2(2\Delta)\rceil}\) and \(K_1=\Delta P\).
Pairwise independence of the colors bounds the expected absolute imbalance of thick and thin outgoing widths in a colony by \((B\sum_i a_i)^{1/2}\), where \(a_i\) is the outgoing core amount of its \(i\)th occurrence. Summing gives a coloring of total imbalance at most \(\sqrt{BQX}\).

For a source colony, group thick outgoing units by their destination pack. There are at most \(K_1\) bins by the length ratio along core edges. Treat the outgoing units of each thin pack as one item and match the entire item into any bin with enough remaining capacity. At the first failure every remaining capacity is below \(2B\), so the unmatched part of the smaller total is at most \(2BK_1\). The one-sided matched width therefore satisfies
\[
 I\ge \tfrac12X-\tfrac12\sqrt{BQX}-2BK_1Q.
\]
The colony bound~\eqref{eq:core-colonies} and the slack condition with \(c_*\ge A\) give \(BQ=O((X+Z+\sigma)/A)\). Take \(A\) a sufficiently large dyadic multiple, up to rounding, of \(1+K_1\). Then \(I\ge X/4-O(Z+\sigma)\) with absolute implicit constant.

Within each bin pair its matched thin units with distinct thick units.
Use Lemma~\ref{lem:powered-gadget} with \(b_*=b\) and \(z=u/b_*\), an
integer by the subdivision hypothesis. In each pair the thick wire plays
\(AB\), with \(A\) its source, and the thin wire plays \(CD\), with \(C\)
its source. The sources share a colony. The destinations \(B,D\) belong
to thick and thin packs, respectively, because a wire stays in one
block and the color is common to that block. Introduce fresh auxiliary
parties \(M,N'\) for each pair. Across all pairs, the two demanded
widths have total \(2I\), and the replacement wires have total width at
most \(2\lambda I\) in each product summand.

In one such summand, add the widths that survive the row contractions
over all pairs. Do the same for the column contractions. Every
replacement wire is internal to one of the two layouts, so the two
totals sum to at most \(2\lambda I\). Perform the cheaper contraction
in every pair. Further identifications can only decrease the surviving
width. The unmatched wires contribute at most \(X-2I\), hence
\[
 X'\le X-(2-\lambda)I\le\rho X+C(Z+\sigma),\qquad
 \rho=1-(2-\lambda)/4<1 .
\]

Matching each thin item in its entirety controls the growth in both load and degree. The merged old packs form stars. A thin pack in its source role joins its own thick pack in the row layout or its one chosen thick destination in the column layout. In its destination role it joins only an auxiliary party. Thus one thick center absorbs at most \(O((1+K_1)A)\) thin packs. Every auxiliary is absorbed into a demand endpoint, to which its new incidences can be charged at a constant multiple of the matched width. A dyadic load multiplier \(g\), polynomial in \(\Delta\), bounds both the resulting core loads and the scaled loads of the frozen lists. The new unit is \(u/b_*\).

Reserve pack slots at every potential dyadic interval. Possible core contacts run between packs of one colony or of intersecting neighboring intervals with ratio at most \(P\). There are \(O(P+\log(2P))\) possible intervals per neighbor for each given interval. The stars and the four matched endpoints have bounded diameter in this graph. Taking a fixed graph power gives a common core support of degree at most \(C_2\Delta^{C_2}\). A tether needs just one old tether adjacency with its own length ratio, followed by a bounded number of core adjacencies. Hence the full degree is at most \(C_2\nabla^{C_2}\), with absolute \(C_2\) enlarged if necessary. There are fewer than \(2r\) potential intervals per position. This gives the claimed bound on positions and keeps tethers attached to their respective frozen tensors.

\subsection{Resetting the core degree}\label{subsec:reset}
Use \(H,N_b,\Delta,\nabla\) for the parameters immediately before the
reset, and put \(P=2^{64\lceil\log_2(2\Delta)\rceil}\),
\(K_1=\Delta P\). The reset replaces the core by a degree-three host.
Its principal cost is a permutation network at each core base. Fix a
trivalent network of size \(c_0Yw_0\) that realizes all permutations
of \(Y\) rows of width \(w_0\) by fixing its tags, where \(4K_1\le Y\)
and \(c_0\ge1\). Call this network the \emph{rotor}.

Figure~\ref{fig:reset} shows the two routes. First partition the blocks
at each position into intervals called \emph{bases}, using their mean
loads, and reserve capacity between each pair of adjacent core bases.
Mixing the blocks makes most bases fit these reservations. A wire
between two such bases crosses a direct link between dedicated rotor
outputs. Every core wire incident to an exceptional base takes a longer
route through binary trees. We must bound the total load on those bases
before charging the extra hops.

\begin{figure}[htb]\centering
\begin{tikzpicture}[x=1cm,y=1cm,
                        font=\footnotesize,term/.style={circle,draw,fill=white,inner sep=2pt}]
% pair of good bases
\foreach \d/\off in {1/-5,-1/5}{
  \begin{scope}[xshift=\off cm,xscale=\d]
    \draw[draw=black!60,rounded corners=6pt](-1.45,-1.5) rectangle (3.36,1.5);
    \node[font=\sffamily\small] at (1,1.8){good base};
    \node at (-.05,1.18){inputs};
    \node at (2.6,1.18){outputs};
    \foreach \y in {.75,0,-.75}{
      \draw (0,\y) -- (2.6,\y);
      \draw (0,\y) -- (-.65,\y) --(-.87,\y+.21);
      \draw (-.65,\y)--(-.87,\y-.21);
    }
    \draw[sea, very thick] (0,0)--(.8,0) (1.8,.75)--(2.6,.75);
    \draw[draw=ink,rounded corners,fill=ink!8] (.8,-1.26) rectangle (1.8,1.12);
    \draw[sea,very thick] (.8,0) -- (1.8,.75);
    \node[font=\scriptsize] at (1.3,-.59){rotor};
    \foreach \y in {.75,0,-.75}{
       \node[term] at (0,\y){};
       \node[term] at (2.6,\y){};
    }
  \end{scope}}
\draw[sea,very thick] (-2.29,.75) -- node[below,text=black]{reserved link} (2.29,.75);

% link from bottom input on left to a terminal in a bad base
\draw[draw=black!60,rounded corners=5pt] (-2.15,-3.94) rectangle (0.2,-2.23);
\node[anchor=west,align=center] at (.42,-3.1) {leaf of a bad base};
\node at (-.61,-2.7) {input};
\draw (-1.41,-3.26) --(-1.68,-2.9);
\draw[rust,thick,densely dashed] (-5.11,-.75)--(-5.65,-.75)--(-5.87,-.96)
  -- (-5.87,-3.6)-- (-1.68,-3.6) -- (-1.41,-3.26)--(-.6,-3.26);
\node[term] at (-.6,-3.26){};
\node[align=center] at (-4,-2.91) {via contact trees};
\end{tikzpicture}
\caption{Routing during a reset. Between good bases, each end of a
wire reaches a dedicated output through a rotor (solid). A contact
with a bad base joins input terminals through binary contact trees
(dashed). Old tensors and their tethers stay at the inputs. Only
selected terminals are shown.}\label{fig:reset}
\end{figure}

\paragraph{Bases and capacities.}
Choose dyadic $F\ge Y^2w_0$; require $c_*\ge F$ in this step. For the mean load \(\mu_v\) at each position over the \(N_b\) blocks, put
\[
 \nu_v=\max\left\{ FB/r,\ 
      \max_{\pi:y\to v} \mu_y\,2^{-\sum_{e\in\pi}s_e/2}\right\},
\]
allowing a path of length zero. The sum of the attenuated weights of paths from each position is bounded by an absolute constant, as in the colony count. Hence \(\sum_v\nu_v\le HFB/r+C\sum_v\mu_v\). At position \(v\) partition into bases of the largest dyadic length \(s_v\le FB/\nu_v\), between \(F\) and \(r\). All partitions of block indices respect their grouping into consecutive $r$'s. Neighboring base depths differ by at most \(s_e/2\). Their total number is \(O((\mathcal M+\sigma)/(FB))\). For each pair of intersecting bases at neighboring core positions \(v,w\), reserve capacity
\[
 \overline y_{vw}\min(s_v,s_w)+FB/K_1 ,
\]
where \(\overline y_{vw}\) is the mean core amount of that contact per block. The first terms sum to at most \(s_v\mu_v\le FB\) at any base, and there are at most \(K_1\) reservations. Hence total reserved capacity at a base is at most \(2FB\). Base lengths and capacities are fixed from the means before permuting the blocks. In each group the output terminals for a given pair of base positions will therefore occupy the same host slots.

\paragraph{Mixing.}
Permute all block indices by one invertible binary matrix. Within the bases, mark dyadic subintervals (including possibly the bases themselves) of load above \(2FB\), propagating with the same edge distances. Leaves have load at most \(2FB\), length at least \(F\), and the previous length-ratio bounds (base splits above the marked region already satisfy propagation by comparison of the base depths; propagate as in colony packing). Call a base bad if it splits or an incident core reservation overflows. For any \(\delta_*>0\), choosing \(F\) sufficiently large in terms of \(\Delta,Y,\delta_*\) gives a permutation with
\begin{equation}
 \sum_{\text{bad bases}}\mathrm{load}\le\delta_*(\mathcal M+\sigma).
 \label{eq:core-badbases}
\end{equation}
To prove this uniformly in the full degree, count how far one overloaded interval can propagate. For a given position having \(n_v=N_b/s_v\) bases, normalize old loads by \(FB\), obtaining \(x_i\le1/F\). At depth \(j\) below the bases let \(a_J\) be the normalized interval loads and \(\overline a\) their mean, at most \(2^{-j}\). Both \(V_j=\sum_J(a_J-\overline a)^2\) and \(U_j=\sum_J a_J^2-\sum_i x_i^2\) bound the number of overloaded intervals and half their (normalized) load. The prefix collision probability for distinct indices under a random invertible binary matrix is at most the reciprocal of the number of intervals. Consequently
\[
 \mathbb E V_j\le\sum_i x_i^2\le n_v/F,\qquad
 \mathbb E U_j\le n_v2^{-j}.
\]
A mark at depth \(j\) can force a base split only along paths with \(\sum s_e\le 2j\), since neighboring base depths change by at most half the edge distance. For this source mark, dyadic nesting gives at most one affected base at each reachable position. The path count used in~\eqref{eq:core-colonies} therefore bounds the number affected by \(2\cdot 2^{j/8}\). At depth \(j\), use \(V_j\) if \(1/F\le2^{-j}\) and \(U_j\) otherwise. Weight this estimator by \(4\cdot 2^{\lceil j/8\rceil}\). Split the sum at \(j=\log_2 F\). The expected charge is at most a constant times \(n_v\sum_{j\ge0}2^{j/8}\min\{F^{-1},2^{-j}\}=O(n_vF^{-7/8})\). Bases with load above \(2FB\) are paid for directly by the depth-zero estimate. For core capacity overflows use the variance estimator for the contact width divided by \(FB\) at the depth of the base intersections. The \(FB/K_1\) buffer above the mean ensures that each overflow contributes at least \(1/K_1^2\). Summing the unweighted variances in expectation over core edges costs at most \(X/(F^2B)\). Each overflow affects at most two bases. Those whose load exceeds \(2FB\) have already been paid for at depth zero, and every other affected base costs at most \(2FB\). Let \(\Phi\) be the sum of the weighted interval estimators and \(K_1^2\) times the contact variances. We have bounded total bad load deterministically by $O(FB\Phi)$ and have shown
\[
 FB\,\mathbb E\Phi = O(F^{-7/8}(\mathcal M+\sigma)+ K_1^2X/F).
\]
A matrix with estimator at most the mean proves \eqref{eq:core-badbases}. Section~\ref{sec:assembly} finds such a permutation by computing collision scores.

\paragraph{The routes.}
After mixing, group consecutive runs of \(r\) blocks. Give every core leaf \(Y\) input terminals. Assign its old tensors whole to a least-loaded terminal successively, so each terminal receives at most \(2FB/Y+B\le 3FB/Y\). Keep all tethers at their assigned terminals, with their other ends in the concatenated frozen list. For every good core base the rotor connects its input terminals to \(Y\) output terminals. A reservation of capacity \(c\) receives \(\lceil cY/(8FB)\rceil\) pairs of dedicated output terminals, joined by direct links. This needs at most \(Y/4+K_1\le Y\) terminals per base. Divide the units of a good--good contact evenly among its pairs. Their output loads are at most \(8FB/Y+u\le9FB/Y\).

At a good base write \(m_{ij}\) for the number of units routed from input \(i\) to output \(j\). The row and column sums of \((\lceil m_{ij}/w_0\rceil)\) are at most
\[
 9FB/(Yuw_0)+Y\le10FB/(Yuw_0).
\]
Bring every row and column sum to the largest sum by adding edges.
Hall's Theorem decomposes the resulting regular bipartite multigraph
into perfect matchings. Each matching configures a layer of the rotor with bit edges powered to width \(u\). Pin unused inputs and corresponding outputs to zero. Take the tensor product of the layer tables at each rotor position.
This uses only the prescribed equality tensor obtained by pinning the rotor tags. Balanced binary trees at each rotor row collect its \(w_0\) ports, at \(O(\log(2w_0))\) hops per unit. Including direct links this costs \(O((c_0+\log(2w_0))FB)\) per good base.

For a core contact with a bad endpoint use balanced binary contact trees attached to its input terminals. The possible neighboring leaf intervals have length ratios at most \(P\), with \(Y\) terminals per leaf, giving depth \(O(\log(2\Delta Y))\). Route ordered lists through the trees by equality tensors. The assigned lists at both ends give a unique internal assignment to the transmitted bits. Choose \(\delta_*\) in \eqref{eq:core-badbases} small enough to pay this routing depth. Thus for an absolute \(A_0\),
\[
 X'\le A_0(1+c_0+\log(2w_0))(X+Z+\sigma).
\]
This counts every core link in the rotors, collecting trees, direct connections and contact trees. The unit remains \(u\). At every input terminal, output terminal, tree vertex, and rotor vertex the load is \(O(FB)\), as follows from the terminal and layer bounds. For a frozen leaf its scaled list load is at most \(2FB\). The core has degree at most three. Rotor vertices have this degree; input terminals have at most rotor-row (via one edge to a row tree) and contact-tree connections within the core; output terminals have at most a rotor-row connection and a direct link. Each binary tree has degree at most three. For tethers the common host allows every input terminal at every possible intersecting core leaf of the prescribed length ratio. The resulting full degree is bounded by a fixed power of \(2\nabla\), with constants depending on this reset. Reserving this hardware at all possible interval slots costs at most \(c r H\) positions.

\subsection{Paying for the reset}
Write
\[
 \kappa=A_0(1+c_0+\log(2w_0))
\]
for the width multiplier of the reset. Take \(t\) to be the least
positive integer such that \(\kappa\rho^t\le1/2\). Repeated use of the
subround estimate gives core width at most
\(\rho^tX+C(Z+\sigma)/(1-\rho)\) just before the reset. After it,
the coefficient of the incoming \(X\) is at most one half. The
remaining cost is a fixed multiple of \(Z+\sigma\), provided we can
choose a rotor with \(Y\ge4K_1\) for the degree reached after \(t\)
subrounds.

Starting at bounded core degree, the recurrence
\(\Delta'\le C_2\Delta^{C_2}\) gives
\[
 \log_2(4K_1)\le C^{t+1}
      \le C''(2+c_0+\log(2w_0))^{C'} .
\]
Here \(C,C',C''\) depend only on the incoming core degree and the
subround constants. For \(Y=2^y\), Theorem~\ref{thm:sort} gives
\(w_0=\lceil y^{a_\zeta}\rceil\) and \(c_0=O_\zeta(y^\zeta)\) for
each \(\zeta>0\). Fix \(\zeta<1/(2C')\), then take \(y\) large. The
displayed polynomial bound is at most \(y\), so this rotor covers the
required core degree.

The number of subrounds \(t\), the rotor, \(F\) and \(\delta_*\)
depend only on the incoming core-degree bound and the conditional
sorting construction. These parameters bound the growth of the maximum
load. The full degree governs tether propagation and the number of host
positions to reserve. All subrounds together divide \(u\) by
\(b_*^t\), and their polynomial degree bounds, followed by the reset,
multiply \(\log(2\nabla)\) by at most a fixed factor.

This completes the proof of Lemma~\ref{lem:core-width}.


\section{The arity budgets after one round}\label{sec:full}

The frozen stations produced by concentration have total arity
$O(L/\sqrt D)$. The core has small maximum arity but may still use more
than $L$ bits. Repeated beats reduce its internal width until the total
charge fits a smaller budget $L/G$. To resume concentration we will
also need a bound on the squared charges and an absolute degree bound
on the full host.

\subsection{From concentration to bounded degree}
Apply Lemma~\ref{lem:local} with arity exponent $\ell=10$. Its absolute degree bound for the core fixes the parameters of a beat in Lemma~\ref{lem:core-width}. Choose
$\epsilon\in(0,1/100]$ so small that in
\[
 J=\lceil2\epsilon\log D\rceil
\]
beats the products of the per-beat increases of the maximum load, of the
reciprocal bundle unit, and of the logarithm of the overall degree are
each at most $D^{0.03}$ when $D$ is sufficiently large. These three
per-beat factors are independent of $D$, which permits this choice.
Increase $D$ below as needed in terms of $k,\delta,\epsilon$.

We require $ed\le L/D$. For each grouping in this round we also require
\begin{equation}\label{eq:round-slack}
  H N_b T/r\le c_D L
\end{equation}
with $c_D>0$ sufficiently small, using the current host counts and the
round's starting $L,T$. The choices of $S,r$ meeting these conditions
in every round are made below.

\paragraph{Concentrating the squared arities.}
Group $r$ blocks and apply the dyadic partition estimate
\eqref{eq:core-colonies} to the charged loads at absolute host degree,
with maximum load $DT/2$ and capacity $DT$. By
\eqref{eq:round-slack} there are $O(L/(DT))$ occupied intervals over
all positions and groups. Regard each interval as one vertex of the graph
to which we apply Lemma~\ref{lem:local}. Include all candidate intervals
as positions, and allow an edge to an interval at an old neighbor only
when the two intersect and satisfy the absolute length-ratio bound of the
dyadic partition. All actual contacts are present. Dyadic nesting bounds
the simple degree of this new host by an absolute constant.

Inside each occupied interval retain each station of charged arity greater
than $T$. Pack the other stations in their order, closing a pack
when its total charged arity reaches $T$. The last pack can be smaller,
and each pack has arity at most $2T$. Its table is the product of its
members, with any internal indices contracted. An interval contains
$O(D)$ resulting stations. Move its external bits, if any, to one
additional identity station, keeping their order and joining the duplicate
ports to the original incidences. Designate these exterior identities as the immobile stations. Their number is at most the number of occupied intervals. The old stations and packs had total arity at most $DT$ per interval, and the identity adds at most $2DT$. The old large stations contribute at most $LT$ to the square sum. Each new pack has arity at most $2T$, so the packs contribute at most $2TL$.
If $f_i$ is the external width of interval $i$, the new identities
contribute at most
\[
 4\sum_i f_i^2\le 4DT\,ed\le 4LT .
\]
With $W=C_1L$ and $C_1$ absolute and sufficiently large, these are
precisely the occupation, immobility, and square bounds of
Lemma~\ref{lem:local}. In particular there are at most $D$ retained large stations, $D+1$ packs and one exterior identity in an occupied interval, within its bound of $100D$.

Expand using Lemma~\ref{lem:local}. Write $Z$ for the total charged
arity of the frozen stations, and $X$ for the total width of the
internal wires between core stations, counted once per wire. In every
summand the frozen charges have total $Z=O(L/\sqrt D)$, square sum
$O(LT)$ and maximum $O(DT)$. The core stations have individual arities
at most $D^{-10}T$ and total arity $O(L(\log D)^{O(1)})$.
The constants in these estimates are uniform in $D$. All positive
internal bundles have width at least $T/D^{O(1)}$.

To apply the packing operations, take an integer unit $u$ below the
minimum guaranteed internal width, still $u\ge T/D^{O(1)}$, divisible
by every subdivision prescribed in the $J$ beats. For fixed $D$ this is
possible once $T$ is large, by rounding down to a multiple of their
product. Pad each positive internal bundle to the next multiple of $u$
with bits pinned to zero. The arity and its square grow by at most
factors two and four. Raise each frozen charge to the larger of its old charge and its padded arity. Its charge grows by at most a factor of two, so the total $Z$ and the square budget still obey the bounds above. Choose a dyadic load bound $B$ of order
$D^{-10}T$ covering every core station. For packing only, scale the
frozen charges by a fixed dyadic multiplier of order $B/(DT)$, small
enough that each is at most $B$. This changes neither the actual widths
nor the frozen square bound. We take $u$ small enough also to ensure $u\le B$.

\paragraph{Removing the core wires.}
Apply Lemma~\ref{lem:core-width} $J$ times, with slack
$\sigma=L/\sqrt D$. The initial full host degree is polynomial in $D$.
The load bounds at intermediate steps are fixed multiples, depending
only on $D$, of $T$, so \eqref{eq:round-slack} supplies the required
slack each time. In the recurrence of Lemma~\ref{lem:core-width} the fixed error $C(Z+\sigma)$ is multiplied by a geometric series of sum less than two. Thus
\[
 X_{\mathrm{end}}\le 2^{-J}X_{\mathrm{start}}+O(Z+\sigma)
      = O\bigl(LD^{-2\epsilon}(\log D)^{O(1)}\bigr).
\]
Frozen stations retain their charges, although one occurrence can now
hold a list of them. Each list has $D^{O(1)}$ members. Indeed, its total scaled load is bounded by the final load bound, while the positive charge floor gives a lower bound on the scaled load of each member. Each core station has
arity at most $D^{-9}T$. The final bundle unit is still at least
$T/D^{O(1)}$, the core degree is at most three, and the overall degree
$\nabla$ satisfies $\log(2\nabla)\le D^{0.04}$ for $D$ sufficiently
large. The latter estimate uses the initial $O(\log D)$ bound on the
logarithm of the degree.

\paragraph{Restoring bounded degree.}
An individual station can meet few \emph{actual} neighbors and still
have many \emph{possible} neighbors in the common host. At a frozen
position reserve a slot for each member of its list. The only frozen contacts are tethers to the core. Route them in three stages (Figure~\ref{fig:tether-trees}). First
split its active tether units into separate outlets, in order. There
are $D^{O(1)}$ outlets, by the maximum width and minimum unit. Next,
from each outlet use a binary tree addressing any potential core
endpoint. At each core endpoint use a binary tree for the potential
incoming outlets. Link corresponding leaves. An internal branching
table is the identity on the concatenated strings of its active
children. Inactive branches need no station. Each block uses the trees
with precisely its own tether assignments, and the trees and slots
themselves are fixed on the common host. The first trees have depth
$O(\log D)$, and the others have depth
$O(\log(2\nabla)+\log D)=O(D^{0.04})$. Their host size increases
by a factor depending only on $D$. A core endpoint retains at most three
core contacts and receives one tree contact, giving overall degree at
most four.

For the linear budget, the total extra arity is $O(ZD^{0.04})$, since each tether bit travels
along one route of that length. For the square budget, the first splitting tree is bounded in terms of the frozen arity. On a fixed level of that tree
the total width is at most the arity of its frozen station. The
sum of squared arities at that level is at most four times the square
of this arity. Summing first over the frozen stations and then over
levels thus gives $O(LT\log D)$. Every other new vertex carries at
most the width of one unit or all the tether bits of one core station,
so has arity at most $2D^{-9}T$. Its contribution to the square sum
is at most $2D^{-9}T$ times its arity. The old core stations have
total arity at most $2X_{\mathrm{end}}+Z$ and obey the same estimate.

\begin{figure}[tb]
\centering
\begin{tikzpicture}[x=1cm,y=.82cm,font=\footnotesize,
  live/.style={draw=sea,line width=1.3pt},
  possible/.style={draw=ink!55,densely dashed},
  branch/.style={circle,draw=ink,fill=white,inner sep=1.35pt},
  station/.style={rounded corners=2pt,draw=ink,inner sep=3pt}]
  \node[station,fill=rust!9] (f) at (0,0) {frozen};
  \node[branch] (split) at (1.2,0) {};
  \node[branch] (a) at (2.9,.78) {};
  \node[branch] (b) at (2.9,-.78) {};
  \node[branch] (c) at (6.45,.78) {};
  \node[branch] (d) at (6.45,-.78) {};
  \node[station,fill=ink!7] (c1) at (7.8,.78) {core};
  \node[station,fill=ink!7] (c2) at (7.8,-.78) {core};
  \draw[live,line width=2.3pt] (f)--(split);
  \draw[live] (split)--(a) (split)--(b);
  \draw[possible] (a)--(3.6,.43)--(5.76,-.43)--(d);
  \draw[possible] (b)--(3.6,-.43)--(5.76,.43)--(c);
  \draw[live] (a)--(3.6,1.12)--(5.76,1.12)--(c)--(c1);
  \draw[live] (b)--(3.6,-1.12)--(5.76,-1.12)--(d)--(c2);
  \foreach \x/\heading in {1.1/{Split into units},3.9/{Select destinations},6.95/{Collect at endpoints}}{
      \node[font=\footnotesize] at (\x,1.75) {\heading};
  }
\end{tikzpicture}
\caption{Two units leave one frozen station. The solid routes are
active in this block. The dashed routes allow either outlet to reach the
other endpoint in a different block. A splitting tree on the left may
carry the full frozen arity for $O(\log D)$ levels. In the middle each
active branch carries a single unit. At the right its width is bounded
by the tether width at the receiving core station.}
\label{fig:tether-trees}
\end{figure}

Together these bounds give charged arities $z_v$ at the end of a
round satisfying
\begin{equation}\label{eq:round-costs}
 \begin{aligned}
 \sum_v z_v&=O\bigl(LD^{-2\epsilon}(\log D)^{O(1)}\bigr),\\
 \sum_v z_v^2&=O\bigl(LT(\log D)^{O(1)}\bigr),\qquad
                   \max_v z_v=O(DT).
 \end{aligned}
\end{equation}
For frozen stations we kept the positive floor from
Lemma~\ref{lem:local}. Every other remaining station meets a positive
internal bundle of width at least $T/D^{O(1)}$. Thus all charges in
\eqref{eq:round-costs} can respect a floor of this order. The constants
and the powers of $\log D$ in \eqref{eq:round-costs} are independent
of $D$. In particular the term $ZD^{0.04}=O(LD^{-0.46})$ fits the
first estimate.

Set $G=2^{\lfloor\epsilon\log D\rfloor}$ and choose a dyadic $q$ with
$G^{1+\delta/4}\le q\le G^{1+\delta/3}$. This interval contains a
dyadic integer for $D$ sufficiently large. Restart with
\begin{equation}\label{eq:round-target}
                 L'=L/G,\qquad T'=qT .
\end{equation}
Indeed $D^{-2\epsilon}(\log D)^{O(1)}=o(1/G)$, and
$L'T'\ge LTG^{\delta/4}$ absorbs the logarithmic loss in the squares.
The maximum in \eqref{eq:round-costs} is at most $DT'/2$ for large $D$.
Finally fix a dyadic $\theta>0$ so small that $\theta qT$ is bounded by
the positive floor just obtained. None of the estimates used in choosing
$G,q,D$ requires a lower bound on $\theta$, since the grouping used the total
load cap and repacked the small stations before applying the local lemma.
This proves \eqref{eq:round-invariants} for the next round.

\subsection{Iteration and the description of the final tables}
\label{iter:subsection}

Fix the round just constructed, including $D$, $G$, $q$ and its $s$
grouping steps. To save a factor $h^a$ in the total charge, take
\[
 j=\left\lceil\frac{a\log h}{\log G}\right\rceil,\qquad
 L_i=L_0G^{-i},\qquad T_i=T_0q^i \qquad(0\le i\le j).
\]
Here $L_i,T_i$ are the budgets after $i$ completed rounds.
It remains to choose $r$ and the initial number $S$ of blocks. A large
$r$ makes partially occupied packs inexpensive. A large $S$ supplies
enough groups for all $sj$ steps and leaves enough blocks to amortize
the common host in Section~\ref{sec:assembly}.

\paragraph{Packing slack and exterior width.}
In one grouping step $N_b$ decreases by $r$ while $H$ increases by at
most $cr$. Their product, which counts allocated occurrences, increases
by at most $c$. The fixed expansions of the host can be included in $c$.
Thus at every stage of the construction,
\begin{equation}
 H N_b\le C^{j+1}mS\log h , \label{iter:occurrences}
\end{equation}
where $C$ depends only on the design of one round. Since
$L_0=C_0 mT_0 S\log(2+h)$, the ratio that controls the slack during
the round starting from $L_i,T_i$ is bounded by
\[
   \frac{H N_b T_i}{rL_i}\le \frac{C_1C^{j+1}(qG)^i}{r}.
\]
Choose $r$ to be the least power of two at least $C_2(CqG)^{j+1}$.
The round design fixes $C_2$ so that this ratio is at most $c_D$ in
\eqref{eq:round-slack} and $r$ meets the lower bound in every step.
In particular $r=h^{O_D(a)+o(1)}$, independently of $S$. The external
width has a separate bound. The $ed$ exposed bits are present in every
round, and $G^j<Gh^a$. Hence for $i\le j$,
\[
       \frac{ed}{L_i}=\frac{eG^i}{C_0m\log(2+h)}
           =O_D\!\left(\frac{h^{-a}}{\log h}\right).
\]
For sufficiently large $h$ this gives $ed\le L_i/D$, the width
allowance for the immobile exterior identities.

\paragraph{Copies and terminal budgets.}
Let $R$ be the least power of two at least $h^M$, for any fixed $M>0$.
Taking $S=Rr^{sj}$ leaves exactly $R$ blocks after $j$ rounds.
Each initial block contains $T_0$ copies, so $d=T_0S$ and
$\log d=O((\log h)^2)$. In particular $d=n^{o(1)}$.
Now $h^a\le G^j<Gh^a$, $q\le G^{1+\delta/3}$ and $T_0=\Theta(h^b)$.
It follows from \eqref{eq:round-start} and \eqref{iter:occurrences} that
\begin{equation}
 \begin{aligned}
 L_j&=O(md\log h\,h^{-a}),&
 T_j&=O\bigl(h^{b+(1+\delta/3)a}\bigr),\\
 HN_b&\le mS h^{O_D(a)+o(1)} .
 \end{aligned}
 \label{eq:round-terminal}
\end{equation}
The occurrence bound holds throughout, with a constant in $O_D(a)$
independent of $M,a,b,\beta$. These scales also justify using a single
working field. They range from $T_0\ge4h^\beta$ to a fixed power of
$h$. Choose the extension degree divisible by the degree of the fixed
extension in Lemma~\ref{lem:local}, and large enough to meet the
cardinality requirement at $T_j$. Its degree is $O(\log h)$ and hence
$O(\log T_i)$ for every $i$.

\paragraph{Descriptions of tables.}
An occurrence carries a label for its table, the widths of its bundles,
and their positions in its port order. The cost of naming its
incident host edges is bounded in terms of $D$ because each edge can
be named by its place in the local neighbor list. Constants indexed by $D$ may depend on the round design. These local data have total length per occurrence at most
\begin{equation}
 C_3^{j+1}(\log h)^{O(1)}
   =h^{O_D(a)+o(1)} .
 \label{iter:labels}
\end{equation}

At initialization the stack at a position has a fixed table label,
specified by its original table, its stack length and the locations
of its ordered sublists. Its description has length
$(\log h)^{O(1)}$. Within one round the following facts bound
the number of old labels that can enter any new label.
\begin{itemize}
\item A colony of capacity $DT_i$ contains at most $D/\theta$ live
entering stations, since each is charged at least $\theta T_i$.
The local expansion in Lemma~\ref{lem:local} uses a bounded number of
old labels, discrete choices and field elements in each output
description. Its additional description length is $O_D(\log h)$.
\item In a subround or reset, both the bound on load per leaf and the
bound on the number of leaves contributing to one resulting pack are
fixed multiples, depending on $D$, of $T_i$ and $1$, respectively.
Each live core contributes at least one internal unit of width
$u\ge T_i/\operatorname{poly}(D)$. The charged load of an individual
frozen station has the same kind of lower bound even after the fixed
rescaling. Scalar stations are withdrawn as soon as they form. It
follows that a contributing pack has a bounded number of old labels
and a bounded number of units. Lists of unit pairings, rotor layers
and ordered sublists at a port therefore have bounded numbers of
entries. Each entry needs at most $O_D(\log h+\log r)$ bits. This
also describes the routing on the final trees at the frozen stations.
\end{itemize}
Thus if $A$ bounds the current label length, one step gives a bound
$C_4A+C_4(\log h)^{C_5}$, including the descriptors of locally
withdrawn scalar factors. The multiplicative constant counts old
labels, independently of the field degree and the length of the input
dictionary labels. Iteration through the $sj$ steps proves
Equation~\eqref{iter:labels}. Global host addresses enter the control
of a whole position only. They never have to be appended to an
individual table description.

\paragraph{The tensor represented by the sum.}
Each term now consists of $R$ surviving blocks and the scalar factors
withdrawn along its history. Regrouping permutes the exterior lists along with their stations and preserves each product entry by entry. A local tensor identity replaces a product by a sum with the same entries on those $ed$ named ports. Colorings, packings, and routings use fixed tie-breaking rules. Only choices from local tensor identities index the expansion.
The final internal wire assignments are summed when the $R$ blocks are
contracted.

For this evaluation one dictionary suffices. A key records the type,
stage, description and arity of a table, followed by its incident bits.
Length fields or fixed padding give a unique encoding. The input
dictionary and the rules of the construction assign a field value
to each valid key. The values at unused keys may be fixed arbitrarily.
In particular, a key for a withdrawn scalar has no incident bits.
A longest key has length
\[
 O_D(T_j)+h^{O_D(a)+o(1)}
\]
by the maximum-arity invariant and Equation~\eqref{iter:labels}.
A derived table description can exceed $(\log h)^{O(1)}$ bits. Section~\ref{sec:assembly} places that description in the \emph{argument} of a dictionary call. The label of the installed table need only specify the key format and length. This accounts separately for the arity of a lookup and the short label required as input to the next macro.

